% संपूर्ण मराठी ग्रंथातील प्रकरण 67: सिद्धता
% हा संपादनयोग्य स्रोतखंड आहे. संकलनासाठी संपूर्ण स्रोतसंग्रहातील
% अक्षररूपे, पूर्वभाग, चिन्हव्याख्या आणि आकृती-स्रोत आवश्यक आहेत.
% मूळ: Open Logic Project; मजकूर आणि रूपांतर CC BY 4.0.
% यंत्रानुवाद, दुरुस्ती आणि तपासणी: OpenAI Codex —
% GPT-5.6 Sol आणि GPT-6 Sol, दोन्ही Ultra effort.
% दुरुस्ती-पुनरावलोकन: GPT-6.1 Sol, Ultra effort.
\chapter{सिद्धता}\label{mth:prf:chap}









\section{प्रस्तावना}\label{mth:prf:int:sec}

प्रास्ताविक तर्कशास्त्राच्या अभ्यासामुळे तुम्हाला एखादी निष्पत्ती-पद्धती परिचित असेल---बहुधा नैसर्गिक निगमनाची किंवा फिच शैलीतील निष्पत्ती-पद्धती, किंवा कदाचित सिद्धता-वृक्षांची पद्धती. या पद्धतींत सूत्र सिद्ध करणे किंवा दिलेला युक्तिवाद युक्त आहे हे दाखवणे, अशा सिद्धता तुम्ही केल्याचे आठवत असेल. त्यासाठी अपेक्षित अंतिम निष्कर्ष मिळेपर्यंत तुम्ही त्या पद्धतीचे नियम वापरले होते. तर्कशास्त्रा\emph{विषयी} विचार करतानाही आपण गोष्टी सिद्ध करतो; पण बहुतेक वेळा निष्पत्ती-पद्धती वापरत नाही. आपण पाहणार असलेल्या बहुतेक सिद्धता संपूर्णपणे प्रथम-क्रम तर्कशास्त्राच्या भाषेत लिहिलेल्या नसून नैसर्गिक भाषेत असतात; या आवृत्तीत त्या मराठीत आहेत आणि त्यांत काही सांकेतिक भाषाही येते. अशा सिद्धता रचताना सुरुवातीला तुम्हाला प्रश्न पडू शकतात: निष्पत्ती-पद्धतीशिवाय एखादी गोष्ट कशी सिद्ध करायची? सुरुवात कुठून करायची? माझी सिद्धता बरोबर आहे हे कसे ठरवायचे?

सिद्धता करण्याचा प्रयत्न सुरू करण्यापूर्वी सिद्धता म्हणजे काय आणि ती कशी रचायची, हे समजणे महत्त्वाचे आहे. नावावरूनच सूचित होते तसे, एखादी गोष्ट सत्य आहे हे दाखवणे हा \emph{सिद्धते}चा उद्देश असतो. संवादाच्या रूपात याचा विचार करता येईल: कोणी तुम्हाला विचारते की एखादी गोष्ट खरी आहे का---उदाहरणार्थ, दोन वगळता प्रत्येक अविभाज्य संख्या विषम आहे का. केवळ “होय” असे उत्तर देणे पुरेसे नाही; त्या व्यक्तीला \emph{का}, हेही जाणून घ्यायचे असेल. अशा वेळी तुम्ही सिद्धता द्याल.

दैनंदिन संभाषणात उत्तराची ढोबळ दिशा सांगणे किंवा अपूर्ण उत्तर देणे पुरेसे ठरू शकते. पण तर्कशास्त्रात आणि गणितात आपल्याला काटेकोर सिद्धता हवी असते: एखादी गोष्ट खरी असल्याविषयी \emph{कोणत्याही} शंकेला जागा राहू नये, असे दाखवायचे असते. याचा अर्थ सिद्धतेतील प्रत्येक पायरीचे समर्थन केले पाहिजे आणि ते समर्थन योग्य असले पाहिजे. उदाहरणार्थ, तुम्ही वापरत असलेले गृहीतक हे सिद्ध करावयाच्या प्रमेयाच्या विधानात खरोखर गृहीत धरलेले असले पाहिजे; व्याख्यांचा वापर अचूक असला पाहिजे; समर्थनासाठी वापरलेली अनुमाने योग्य असली पाहिजेत; इत्यादी.

सहसा आपण एखादे विधान सिद्ध करत असतो. सिद्ध करावयाच्या विधानांना आपण विविध नावे देतो: विधान, प्रमेय, पूर्वप्रमेय किंवा अनुप्रमेय. विधान म्हणजे सिद्ध करावे असे मूलभूत म्हणणे: नोंदवण्याइतके महत्त्वाचे, पण कदाचित फार सखोल नसलेले किंवा वारंवार न वापरले जाणारे. प्रमेय म्हणजे लक्षणीय, महत्त्वाचे विधान. त्याची सिद्धता अनेकदा काही पायऱ्यांत विभागलेली असते. कधी त्याला ते प्रथम सिद्ध करणाऱ्या व्यक्तीचे नाव दिले जाते---उदाहरणार्थ, कँटरचे प्रमेय किंवा लोव्हेनहाइम--स्कोलेम प्रमेय (L\"owenheim--Skolem)---तर कधी त्याच्या विषयावरून नाव दिले जाते---उदाहरणार्थ, संपूर्णता प्रमेय. पूर्वप्रमेय म्हणजे अधिक महत्त्वाचा निष्कर्ष सिद्ध करण्यासाठी वापरले जाणारे विधान किंवा प्रमेय. नावामुळे गोंधळ होऊ शकतो: कधी पूर्वप्रमेये स्वतःही महत्त्वाचे निष्कर्ष असतात आणि ती मांडणाऱ्या व्यक्तीच्या नावाने ओळखली जातात---उदाहरणार्थ, झोर्नचे पूर्वप्रमेय. अनुप्रमेय म्हणजे दुसऱ्या निष्कर्षावरून सहज मिळणारा निष्कर्ष.

सिद्ध करावयाच्या विधानात अनेकदा काही गृहीतके असतात. आपण कोणत्या प्रकारच्या गोष्टींविषयी सिद्धता करतो, हे त्यांमुळे स्पष्ट होते. विधानाची सुरुवात “$\varphi$ हे $\psi \lif \chi$ या रूपाचे सूत्र असू द्या” किंवा “$\Gamma \Proves \varphi$ असे समजा”, किंवा अशाच शब्दांनी होऊ शकते. ही त्या विधानाची, प्रमेयाची किंवा पूर्वप्रमेयाची \emph{गृहीतके} आहेत; सिद्धतेत ती सत्य आहेत असे धरता येते. काय सिद्ध करायचे आहे याची व्याप्ती ती निश्चित करतात आणि आपण ज्या वस्तूंविषयी बोलतो त्यांच्यासाठी काही नावेही देतात. उदाहरणार्थ, तुमच्या विधानाची सुरुवात “$\varphi$ हे $\psi \lif \chi$ या रूपाचे सूत्र असू द्या” अशी असेल, तर तुम्ही केवळ एका विशिष्ट प्रकारच्या सर्व सूत्रांविषयी---म्हणजे सशर्त सूत्रांविषयी---काहीतरी सिद्ध करत आहात. तुमची सिद्धता ज्या $\psi \lif \chi$ विषयी बोलते ते त्या प्रकारचे स्वेच्छ सशर्त सूत्र आहे, असे समजले जाते.












\section{सिद्धतेची सुरुवात}\label{mth:prf:str:sec}

पण सुरुवात करायची तरी कुठून?

तुम्हाला काहीतरी सिद्ध करायला दिले आहे. म्हणून सिद्धतेत शेवटी त्याचाच उल्लेख आला पाहिजे. अर्थात, सुरुवातीलाच आपण ते सिद्ध करणार आहोत असे \emph{जाहीर} करता येते; पण ते गृहीतक म्हणून वापरायचे नाही. कोऱ्या कागदाच्या तळाशी काय सिद्ध करायचे आहे ते लिहा. त्यामुळे तुमचे ध्येय नजरेआड होणार नाही.

यानंतर, वापरता येण्याजोगी काही गृहीतके असतील. तुम्ही करत असलेल्या सिद्धतेच्या \emph{प्रकारा}विषयी पुढील विभागांत बोलताना हे अधिक स्पष्ट होईल. ही गृहीतके पानाच्या वरच्या बाजूला लिहा आणि ती गृहीतके आहेत हे स्पष्टपणे दर्शवा. उदाहरणार्थ, $p$ गृहीत धरत असाल, तर “$p$ असे गृहीत धरा” किंवा “$p$ असे समजा” असे लिहा. शेवटी, प्रश्नात काही व्याख्या असतील ज्या तुम्हाला समजणे आवश्यक आहे. एखादी विशिष्ट व्याख्या वापरायला सांगितलेली असेल; किंवा गृहीतकांत आणि ज्या निष्कर्षापर्यंत पोहोचायचे आहे त्यात विविध व्याख्या आलेल्या असतील. \emph{या व्याख्या लिहून काढा आणि त्यांचा अर्थ तुम्हाला समजतो आहे याची खात्री करा.}

सिद्धतेची मांडणी कशी करायची हे प्रश्नाच्या रूपावरही अवलंबून असते. वाक्याच्या प्रकारानुसार सिद्धतेची मांडणी कशी करायची याचा तपशील पुढील विभागांत दिला आहे.












\section{व्याख्यांचा वापर}\label{mth:prf:def:sec}

सिद्धतेत वापरल्या जाऊ शकणाऱ्या सर्व व्याख्या तुम्हाला परिचित असल्या पाहिजेत आणि त्यांचा योग्य वापर करता आला पाहिजे, असे आपण म्हटले होते. हा मुद्दा फार महत्त्वाचा आहे; त्याचा थोडा अधिक तपशील पाहू. गुणधर्म आणि संबंधांविषयी थोडक्यात बोलता यावे म्हणून व्याख्यांद्वारे त्यांना संक्षिप्त नावे दिली जातात. असे नाव म्हणजे \emph{व्याख्येय}; ते ज्याचा संक्षेप करते तो मजकूर म्हणजे \emph{व्याख्यापक}. सिद्धतेत अनेकदा व्याख्येय कसे मांडले होते याकडे परत जावे लागते. कारण सिद्धता पुढे नेण्यासाठी व्याख्यापकाची तार्किक रचना वापरावी लागते: व्याख्येय संज्ञा ज्याचा संक्षेप आहे तोच हा विस्तृत मजकूर असतो. व्याख्या उलगडल्यामुळे प्रत्यक्ष तार्किक कार्य जिथे होते तिथपर्यंत आपण पोहोचतो.

एका उदाहरणाने सुरुवात करू. पुढील विधान सिद्ध करायचे आहे असे समजा:

\begin{prop}
कोणत्याही $A$ आणि $B$ संचांसाठी, $A \cup B = B \cup A$.
\end{prop}

सिद्धतेची सुरुवात करण्यासाठीच दोन संच एकरूप असणे म्हणजे काय हे समजणे आवश्यक आहे. म्हणजे त्या समीकरणातील “$=$” चा संचांसाठी काय अर्थ आहे हे समजले पाहिजे. दोन संचांचे घटक तेच असतील तेव्हा ते संच एकरूप असतात, अशी व्याख्या आहे. म्हणून उलगडायची व्याख्या ही आहे:

\begin{defn}
$A$ आणि $B$ संच \emph{एकरूप} आहेत, $A = B$, तेव्हाच आणि फक्त तेव्हा जेव्हा~$A$ मधील प्रत्येक घटक हा~$B$ मधील घटक असतो आणि उलटही तसेच असते.
\end{defn}

या व्याख्येत $A$ आणि~$B$ ही स्वेच्छ संचांसाठी वापरलेली प्रतिनिधी नावे आहेत. ही व्याख्या ज्या गोष्टीची व्याख्या करते---म्हणजे \emph{व्याख्येय}---ती “$A = B$” ही अभिव्यक्ती आहे. $A = B$ कोणत्या अटीवर सत्य असते हे व्याख्या सांगते. ती अट---“~$A$ मधील प्रत्येक घटक हा~$B$ मधील घटक असतो आणि उलटही तसेच असते”---म्हणजे \emph{व्याख्यापक}.\footnote{या विशिष्ट प्रकरणात---आणि येथे फार गोंधळ होऊ शकतो!---$A = B$ असते तेव्हा $A$ आणि $B$ हे प्रत्यक्षात एकच संच असतात, जरी त्यासाठी डावीकडे आणि उजवीकडे वेगवेगळी अक्षरे वापरली असली तरी. पण त्या संचाची ओळख करून देण्याचे मार्ग वेगळे असू शकतात; त्यामुळे व्याख्या निरर्थक पुनरुक्ती ठरत नाही.} ही अट पूर्ण होते तेव्हाच आणि फक्त तेव्हा (इंग्रजी संक्षेप “iff”) $A = B$ सत्य असते, असे ही व्याख्या सांगते.

व्याख्या वापरताना तिच्यातील $A$ आणि $B$ यांचा संबंध प्रत्यक्ष हाताळत असलेल्या प्रकरणाशी योग्य रीतीने जुळवला पाहिजे. आपल्या प्रकरणात $A \cup B = B \cup A$ सत्य असण्यासाठी प्रत्येक $z \in A \cup B$ हा $B \cup A$ मध्येही असला पाहिजे आणि उलटही तसेच असले पाहिजे. विधानातील $A \cup B$ ही अभिव्यक्ती व्याख्येतील~$A$ ची भूमिका बजावते आणि $B \cup A$ ही व्याख्येतील~$B$ ची. व्याख्येत आणि आपण सिद्ध करत असलेल्या विधानात $A$ आणि $B$ हीच नावे वेगवेगळ्या भूमिकांत येतात; म्हणून या दोन वापरांचा गोंधळ होणार नाही याची काळजी घ्यावी. उदाहरणार्थ,~$A$ मधील प्रत्येक घटक हा~$B$ मधील घटक आहे आणि उलटही तसेच आहे, हे दाखवून विधान सिद्ध करता येईल असे समजणे चुकीचे होईल. त्यामुळे $A = B$ सिद्ध होईल; $A \cup B = B \cup A$ सिद्ध होणार नाही. शिवाय, $A$ आणि $B$ हे कोणतेही दोन संच असू शकतात. $A$ आणि~$B$ यांच्याविषयी काहीच गृहीत धरलेले नसेल, तर ते वेगळे संचही असू शकतात; म्हणून त्या चुकीच्या मार्गाने फार पुढे जाता येणार नाही.

सिद्धतेत संयोगासारख्या संचसैद्धान्तिक संकल्पना येतात. त्यामुळे सिद्धता कशी पुढे न्यायची हे समजण्यासाठी $\cup$ या चिन्हाचाही अर्थ समजला पाहिजे. कधी एखादी व्याख्या उलगडल्यावर आणखी व्याख्या उलगडाव्या लागतात. उदाहरणार्थ, $A \cup B$ ची व्याख्या $\Setabs{z}{z \in A \text{ किंवा } z \in B}$ अशी आहे. म्हणून $x \in A \cup B$ सिद्ध करायचे असेल, तर $\cup$ ची व्याख्या उलगडल्यावर $x \in \Setabs{z}{z \in A \text{ किंवा } z \in B}$ सिद्ध करावे लागेल असे दिसते. आता $x \in \Setabs{z}{\dots z\dots}$ हे $\dots x\dots$ असते तेव्हाच आणि फक्त तेव्हाच खरे असते, हेही आठवले पाहिजे. म्हणून $\Setabs{z}{\dots z \dots}$ या लेखनाची व्याख्या आणखी उलगडल्यावर असे दिसते की $x \in A$ किंवा $x \in B$ हे दाखवायचे आहे. त्यामुळे “$A \cup B$ मधील प्रत्येक घटक हा $B \cup A$ मधील घटक देखील असतो” याचा प्रत्यक्ष अर्थ असा आहे: “प्रत्येक $x$ साठी, $x \in A$ किंवा $x \in B$ असेल, तर $x \in B$ किंवा $x \in A$ असते.” विधानातील व्याख्या पूर्ण उलगडल्यावर आपण हे दाखवायचे आहे असे दिसते:

\begin{prop}
कोणत्याही $A$ आणि $B$ संचांसाठी: (a) प्रत्येक $x$ साठी, $x \in A$ किंवा $x \in B$ असेल, तर $x \in B$ किंवा $x \in A$ असते; आणि (b) प्रत्येक $x$ साठी, $x \in B$ किंवा $x \in A$ असेल, तर $x \in A$ किंवा $x \in B$ असते.
\end{prop}

व्याख्या उलगडणे हा सिद्धता रचण्याचा आवश्यक भाग आहे, हा महत्त्वाचा मुद्दा आहे. हे योग्य रीतीने करणे कधी अवघड असते: व्याख्येतील चल आणि सिद्ध करावयाच्या दाव्यातील पदे यांतील फरक ओळखून त्यांची योग्य जुळवणी केली पाहिजे. यशस्वी होण्यासाठी प्रश्न नेमके काय विचारतो आणि त्यातील सर्व संज्ञांचा अर्थ काय आहे हे समजले पाहिजे. अनेकदा एकाहून अधिक व्याख्या उलगडाव्या लागतील. पुढे येणाऱ्या साध्या सिद्धतांत उत्तर जवळजवळ थेट व्याख्यांवरूनच मिळते. अर्थात, प्रत्येक वेळी हे इतके सोपे असेलच असे नाही.

\begin{prob}
$A \cap B \neq \emptyset$ सिद्ध करायला सांगितले आहे असे समजा. येथे येणाऱ्या सर्व व्याख्या उलगडा; म्हणजे “$\cap$”, “=” किंवा “$\emptyset$” यांचा उल्लेख न करता हे विधान पुन्हा मांडा.
\end{prob}












\section{अनुमानाचे नमुने}\label{mth:prf:pat:sec}

सिद्धता स्वतंत्र अनुमानांनी बनतात. अनुमान करताना आपण सहसा “म्हणून”, “त्यामुळे” किंवा “यावरून” असे शब्द वापरून ते दर्शवतो. त्या अनुमानाचा आधार अनेकदा सिद्धतेत आधीच उपलब्ध असलेल्या एक-दोन गोष्टी असतात: एखादे गृहीतक किंवा आधीच्या अनुमानाने मिळालेला निष्कर्ष. स्पष्टतेसाठी या गोष्टींना नावे किंवा क्रमांक देता येतात आणि अनुमानात कोणती इतर विधाने वापरली आहेत हे दर्शवता येते. नव्या निष्कर्षाची सत्यता आधीच्या गोष्टींवरून \emph{का} मिळते याचे स्पष्टीकरणही अनुमानात अनेकदा असते. सिद्धतांत वारंवार वापरले जाणारे अनुमानाचे काही सामान्य नमुने आहेत; त्यांपैकी काही पुढे पाहू. विगमनाने सिद्धता करण्यासारखे काही नमुने अधिक गुंतागुंतीचे आहेत; त्यांची चर्चा नंतर होईल.

अनुमानाचा एक नमुना आपण आधीच पाहिला आहे: व्याख्या उलगडणे किंवा तिचा वापर करणे. व्याख्या उलगडताना व्याख्येय असलेले म्हणणे आपण व्याख्यापकाच्या रूपात पुन्हा मांडतो. उदाहरणार्थ, सिद्धतेत आधीच $D = E$ हे (a) प्रस्थापित केले आहे असे समजा. मग संचांसाठीच्या $=$ च्या व्याख्येवरून असे अनुमान करता येते: “म्हणून (a) आणि व्याख्येवरून,~$D$ मधील प्रत्येक घटक हा~$E$ मधील घटक असतो आणि उलटही तसेच असते.”

थोडा गोंधळ होऊ शकतो: अनेकदा अनुमानाचे समर्थन प्रत्यक्ष अनुमान करताना लिहिण्याऐवजी त्याआधीच लिहिले जाते. $D = E$ अद्याप सिद्ध केलेले नाही, पण ते सिद्ध करायचे आहे असे समजा. अपेक्षित निष्कर्ष $D = E$ असेल, तर व्याख्या वापरून हे ध्येयही पुन्हा मांडता येते: $D = E$ सिद्ध करण्यासाठी~$D$ मधील प्रत्येक घटक हा~$E$ मधील घटक आहे आणि उलटही तसेच आहे, हे सिद्ध करावे लागते. म्हणून सिद्धतेचे रूप असे असेल: (a)~$D$ मधील प्रत्येक घटक हा~$E$ मधील घटक आहे हे सिद्ध करा; (b)~$E$ मधील प्रत्येक घटक हा~$D$ मधील घटक आहे हे सिद्ध करा; (c) म्हणून (a), (b) आणि $=$ च्या व्याख्येवरून $D = E$. पण सहसा आपण ते अशा रीतीने लिहीत नाही. त्याऐवजी असे लिहू शकतो:
\begin{quote}
$D = E$ दाखवायचे आहे. ~$=$ च्या व्याख्येनुसार त्यासाठी~$D$ मधील प्रत्येक घटक हा~$E$ मधील घटक आहे आणि उलटही तसेच आहे, हे दाखवणे पुरेसे आहे.

(a) \dots (~$D$ मधील प्रत्येक घटक हा~$E$ मधील घटक आहे याची सिद्धता) \dots

(b) \dots (~$E$ मधील प्रत्येक घटक हा~$D$ मधील घटक आहे याची सिद्धता) \dots
\end{quote}

\subsection{संधियोगाचा वापर}

संधियोगातील एखादा संधिघटक निष्कर्ष म्हणून काढणे हा कदाचित अनुमानाचा सर्वात साधा नमुना आहे. म्हणजे $p$ आणि~$q$ असे गृहीत धरले असेल किंवा आधीच सिद्ध केले असेल, तर~$p$ चे (आणि~$q$ चेही) अनुमान करता येते. हे इतके मूलभूत अनुमान आहे की अनेकदा त्याचा स्वतंत्र उल्लेख होत नाही. उदाहरणार्थ, $D = E$ ची व्याख्या उलगडल्यावर~$D$ मधील प्रत्येक घटक हा~$E$ मधील घटक आहे आणि उलटही तसेच आहे, हे आपल्याला मिळते. यावरून~$E$ मधील प्रत्येक घटक हा~$D$ मधील घटक आहे असा निष्कर्ष काढता येतो; तोच “उलटही तसेच” हा भाग आहे.

\subsection{संधियोग सिद्ध करणे}

कधी सिद्ध करावयाचे विधान संधियोगाच्या रूपात असेल: “$p$ आणि $q$ सिद्ध करा” असे सांगितले जाईल. अशा वेळी दोन गोष्टी कराव्या लागतात: $p$ सिद्ध करा आणि नंतर $q$ सिद्ध करा. सिद्धता दोन भागांत विभागून स्पष्टतेसाठी त्या भागांना नावे किंवा क्रमांक देता येतात. सुरुवातीच्या टिपणांत पानाच्या वरच्या बाजूला “(1) $p$ सिद्ध करा” आणि मध्यभागी “(2) $q$ सिद्ध करा” असे लिहिता येईल. अर्थात, संधियोग सिद्ध करायला स्पष्टपणे सांगितलेले नसले तरी सिद्धतेत संधियोग सिद्ध करणे आवश्यक ठरू शकते. उदाहरणार्थ, $D = E$ सिद्ध करायला सांगितले असेल, तर~$=$ ची व्याख्या उलगडल्यावर~$D$ मधील प्रत्येक घटक हा~$E$ मधील घटक आहे \emph{आणि}~$E$ मधील प्रत्येक घटक हा~$D$ मधील घटक आहे, हे सिद्ध करावे लागते असे दिसेल.

\subsection{विकल्पयोग सिद्ध करणे}

सिद्ध करावयाचे विधान विकल्पयोगाच्या रूपात असेल---म्हणजे “$p$ किंवा $q$” या रूपाचे असेल---तर एका विकल्पघटकाची सत्यता दाखवणे पुरेसे असते. पण प्रमेयाच्या गृहीतकांवरून कोणताही एक विकल्पघटक आपोआप मिळतो, असे सहसा घडत नाही. उलट, अनेकदा प्रमेयाची गृहीतकेच विकल्पयोगात्मक असतात; किंवा एका प्रकारच्या सर्व वस्तूंत दोनपैकी एक गुणधर्म आहे हे दाखवायचे असते, पण काही वस्तूंत पहिला आणि इतरांत दुसरा गुणधर्म असतो. अशा वेळी प्रकरणांनुसार सिद्धता उपयुक्त ठरते; ती पुढे पाहू.

\subsection{सोपाधिक सिद्धता}

अनेक प्रमेये सशर्त रूपात असतात: $p$ असेल, तर $q$ देखील सत्य आहे, हे दाखवायचे असते. अशा सिद्धतेची मांडणी सहज करता येते: सशर्त विधानाचे पूर्वांग---येथे $p$---गृहीत धरा आणि त्यावरून निष्कर्ष~$q$ सिद्ध करा. म्हणून प्रमेय “$p$ असेल, तर $q$” असे असेल, तर सिद्धतेची सुरुवात “$p$ असे गृहीत धरा” अशी करा आणि शेवटी~$q$ सिद्ध झालेले असले पाहिजे.

सशर्त विधाने वेगवेगळ्या रीतीने मांडता येतात. “$p$ असेल, तर $q$” याऐवजी प्रमेय “$q$ असेल तरच $p$”, “$p$ असेल तर $q$” किंवा “$p$ या अटीवर $q$” असे म्हणू शकते. या सर्वांचा एकच अर्थ आहे: $p$ गृहीत धरून त्यावरून~$q$ सिद्ध करावे लागते. द्विसशर्त विधान---“$p$ सत्य असते तेव्हाच आणि फक्त तेव्हा जेव्हा $q$ सत्य असते” (इंग्रजी संक्षेप “iff”)---हे दोन सशर्त विधानांचे एकत्रीकरण असते, हे आठवा: $p$ असेल तर $q$ आणि $q$ असेल तर~$p$. म्हणून सोपाधिक सिद्धता दोनदा करावी लागते: एकदा पहिल्या सशर्त विधानासाठी आणि एकदा दुसऱ्यासाठी. मात्र कधी इतर “तेव्हाच आणि फक्त तेव्हा” विधानांची साखळी जोडूनही असे विधान सिद्ध करता येते: सुरुवात “$p$” पासून होते आणि शेवट “$q$” येथे. अशा वेळी प्रत्येक पायरी खरोखर “तेव्हाच आणि फक्त तेव्हा” अशीच आहे याची खात्री केली पाहिजे.

\subsection{वैश्विक विधाने}

वैश्विक विधानाचा वापर सोपा आहे: एखादी गोष्ट प्रत्येक वस्तूसाठी खरी असेल, तर ती प्रत्येक विशिष्ट वस्तूसाठी खरी असते. उदाहरणार्थ, सिद्धतेचे गृहीतक $A \subseteq B$ असेल, तर~$\subseteq$ ची व्याख्या उलगडल्यावर त्याचा अर्थ असा होतो: प्रत्येक $x \in A$ साठी $x \in B$ असते. त्यामुळे $z \in A$ हे आधीच माहीत असेल, तर $z \in B$ असा निष्कर्ष काढता येतो.

वैश्विक विधान सिद्ध करणे थोडे अवघड वाटू शकते. सहसा ही विधाने “$x$ मध्ये~$P$ हा गुणधर्म असेल, तर त्यात~$Q$ हाही गुणधर्म असतो” किंवा “सर्व $P$ असलेल्या वस्तू $Q$ असलेल्या असतात” या रूपात असतात. अर्थात, प्रत्येक विधान नेमके या रूपात असेलच असे नाही; नेमके काय सिद्ध करायला सांगितले आहे हे ओळखण्यासाठी थोडा सराव लागतो. पण अनेकदा एखादा गुणधर्म असलेल्या सर्व वस्तूंमध्ये आणखी एक विशिष्ट गुणधर्म आहे, हे सिद्ध करावे लागते.

वैश्विक विधान सिद्ध करण्यासाठी पहिला गुणधर्म असलेल्या वस्तूंना नावे किंवा चल द्यायचे आणि त्यांत दुसराही गुणधर्म आहे हे दाखवायचे. हे असेही म्हणता येईल: \emph{सर्व}~$P$ असलेल्या वस्तूंसाठी काहीतरी सिद्ध करायचे असेल, तर एका \emph{स्वेच्छ}~$P$ असलेल्या वस्तूसाठी ते सिद्ध करावे लागते. दिलेले नाव हे अशा स्वेच्छ~$P$ वस्तूचे नाव असते. स्वेच्छ वस्तूंसाठी सहसा एकेक अक्षर वापरतो आणि त्या अक्षरांसाठी काही रूढी असतात: नैसर्गिक संख्यांसाठी $n$, सूत्रे साठी $\varphi$, संचांसाठी $A$, फलनांसाठी $f$, इत्यादी.

$P$ या गुणधर्मापलीकडे विशेष अटी न घालता संपूर्ण सिद्धतेत सामान्यता टिकवणे हा मुख्य मुद्दा आहे. एका स्वेच्छ वस्तूत (“$x$”)~$P$ हा गुणधर्म आहे असे गृहीत धरून सुरुवात करायची आणि केवळ व्याख्या किंवा अनुमत गृहीतकांच्या आधारे $x$ मध्ये~$Q$ हा गुणधर्म आहे हे दाखवायचे. $x$ नेमके काय आहे याविषयी आणखी कोणतीही विशेष अट घातलेली नसल्यामुळे, $P$ असलेल्या प्रत्येक वस्तूत~$Q$ आहे असे म्हणता येते. म्हणजे $x$ हे~$P$ गुणधर्म असलेल्या \emph{सर्व} वस्तूंचे प्रतिनिधी नाव आहे.

\begin{prop}
सर्व $A$ आणि $B$ संचांसाठी, $A \subseteq A \cup B$.
\end{prop}

\begin{proof}
$A$ आणि $B$ हे स्वेच्छ संच असू द्या. $A \subseteq A \cup B$ दाखवायचे आहे. $\subseteq$ च्या व्याख्येनुसार याचा अर्थ असा: प्रत्येक $x$ साठी, $x \in A$ असेल तर $x \in A \cup B$ असते. म्हणून $x \in A$ हा~$A$ मधील स्वेच्छ घटक असू द्या. $x \in A \cup B$ दाखवायचे आहे. $x \in A$ असल्यामुळे $x \in A$ किंवा $x \in B$ असते. म्हणून $x \in \Setabs{x}{x \in A \lor x \in B}$. पण $\cup$ च्या व्याख्येनुसार याचाच अर्थ $x \in A \cup B$ असा होतो.
\end{proof}

\subsection{प्रकरणांनुसार सिद्धता}

गृहीतक म्हणून किंवा आधीच मिळालेला निष्कर्ष म्हणून विकल्पयोग उपलब्ध आहे असे समजा: $p$ किंवा $q$ सत्य आहे असे गृहीत धरलेले आहे किंवा सिद्ध झालेले आहे. तुम्हाला $r$ सिद्ध करायचे आहे. हे दोन पायऱ्यांत करता येते: प्रथम $p$ सत्य आहे असे गृहीत धरून~$r$ सिद्ध करा; नंतर $q$ सत्य आहे असे गृहीत धरून पुन्हा~$r$ सिद्ध करा. दोन पर्यायांपैकी एक खरा आहे असे गृहीत धरलेले किंवा माहीत असल्यामुळे ही पद्धत चालते. या दोन पायऱ्यांमुळे कोणताही एक पर्याय~$r$ च्या सत्यतेसाठी पुरेसा आहे हे मिळते. दोन्ही पर्याय खरे असतील, तर $r$~खरे असण्याची एकाऐवजी दोन कारणे असतील. $p$ आणि~$q$ दोन्ही गृहीत धरून $r$~सत्य आहे हे वेगळे सिद्ध करण्याची गरज नाही. काय करत आहोत हे दर्शवण्यासाठी “प्रकरणे वेगळी पाहू” असे जाहीर करतो. उदाहरणार्थ, $x \in B \cup C$ माहीत आहे असे समजा. $B \cup C$ ची व्याख्या $\Setabs{x}{x \in B \text{ किंवा } x \in C}$ अशी आहे. म्हणजे व्याख्येनुसार $x \in B$ किंवा $x \in C$. यावरून $x \in A$ सिद्ध करण्यासाठी प्रथम $x \in B$ गृहीत धरून त्यावरून $x \in A$ सिद्ध करू; नंतर $x \in C$ गृहीत धरून पुन्हा $x \in A$ सिद्ध करू. गृहीतकाखाली “प्रकरणे वेगळी पाहू” लिहा; त्याखाली “प्रकरण (1): $x \in B$” लिहा आणि पानाच्या मध्यावर “प्रकरण (2): $x \in C$” लिहा. मग वरच्या आणि खालच्या भागांतील सिद्धता पूर्ण करा.

सिद्ध करावयाचे विधान स्वतः विकल्पयोगात्मक असेल, तर प्रकरणांनुसार सिद्धता विशेष उपयुक्त ठरते. हे एक साधे उदाहरण:

\begin{prop}
$B \subseteq D$ आणि $C \subseteq E$ असे समजा. मग $B \cup C \subseteq D \cup E$.
\end{prop}

\begin{proof}
(a) $B \subseteq D$ आणि (b) $C \subseteq E$ असे गृहीत धरा. व्याख्येनुसार, कोणताही $x \in B$ हा $\in D$ देखील असतो (c) आणि कोणताही $x \in C$ हा $\in E$ देखील असतो (d). $B \cup C \subseteq D \cup E$ दाखवण्यासाठी $x \in B \cup C$ असेल, तर $x \in D \cup E$ असते हे दाखवावे लागते; हे $\subseteq$ च्या व्याख्येवरून मिळते. ~$ \cup$ च्या व्याख्येनुसार $x \in B \cup C$ हे $x \in B$ किंवा $x \in C$ असते तेव्हाच आणि फक्त तेव्हाच सत्य असते. त्याचप्रमाणे $x \in D \cup E$ हे $x \in D$ किंवा $x \in E$ असते तेव्हाच आणि फक्त तेव्हाच सत्य असते. म्हणून हे दाखवावे लागते: कोणत्याही $x$ साठी, $x \in B$ किंवा $x \in C$ असेल, तर $x \in D$ किंवा $x \in E$ असते.

\begin{quote}
आत्तापर्यंत फक्त व्याख्या उलगडल्या आहेत!{} आपल्या विधानाची $\subseteq$ आणि $\cup$ न वापरता पुनर्मांडणी केली आहे; आता एक वैश्विक सशर्त विधान सिद्ध करायचे उरले आहे. आधीच्या चर्चेनुसार, $x$ ही अशी वस्तू आहे की तिच्याविषयी “जर” भाग सत्य आहे असे गृहीत धरायचे आणि मग “तर” भागही सत्य आहे हे दाखवायचे. म्हणजे $x \in B$ किंवा $x \in C$ असे गृहीत धरून $x \in D$ किंवा $x \in E$ दाखवू.\footnote{हा परिच्छेद आपण काय करत आहोत याचे स्पष्टीकरण देतो; तो सिद्धतेचा भाग नाही. स्वतःच्या सिद्धता लिहिताना हा सगळा तपशील देण्याची गरज नाही.}
\end{quote}

$x \in B$ किंवा $x \in C$ असे समजा. $x \in D$ किंवा $x \in E$ दाखवायचे आहे. प्रकरणे वेगळी पाहू.

प्रकरण 1: $x \in B$. (c) वरून $x \in D$. म्हणून $x \in D$ किंवा $x \in E$. येथे एका विकल्पघटकावरून विकल्पयोग मिळवण्याचे आधी चर्चिलेले अनुमान केले आहे.

प्रकरण 2: $x \in C$. (d) वरून $x \in E$. म्हणून $x \in D$ किंवा $x \in E$.
\end{proof}

\subsection{अस्तित्ववाची विधान सिद्ध करणे}

अस्तित्ववाची विधान सिद्ध करायला सांगितले असेल, तर प्रश्नाचे रूप सहसा “असा~$x$ आहे की $\dots x \dots$, हे सिद्ध करा” असे असते. म्हणजे “$\dots x \dots$” ने वर्णन केलेला गुणधर्म असलेली वस्तू अस्तित्वात आहे, हे दाखवायचे असते. अशा वेळी योग्य वस्तू ओळखून तिच्यात आवश्यक गुणधर्म आहे हे दाखवावे लागते. हे सरळ वाटते; पण अशा सिद्धता अवघडही असू शकतात. सहसा त्यात एखादी वस्तू \emph{रचावी} किंवा \emph{परिभाषित} करावी लागते आणि अशा रीतीने परिभाषित केलेल्या वस्तूत आवश्यक गुणधर्म आहे हे सिद्ध करावे लागते. योग्य वस्तू शोधणे अवघड असू शकते; तिच्यात आवश्यक गुणधर्म आहे हे सिद्ध करणे अवघड असू शकते; आणि कधी आपण मुळात एखादी वस्तू परिभाषित करण्यात यशस्वी झालो आहोत हेच दाखवणे अवघड असते!{}

सहसा वस्तू स्पष्टपणे सांगून अशी सिद्धता लिहितो. उदाहरणार्थ, “$x$ म्हणजे \dots{} असू द्या” असे लिहितो; येथे \dots{} ने कोणती वस्तू अभिप्रेत आहे ते सांगितले जाते. आवश्यक असल्यास $\dots$ खरोखर अस्तित्वात असलेल्या वस्तूचे वर्णन करते हे सिद्ध करतो; मग $x$ मध्ये~$Q$ हा गुणधर्म आहे हे दाखवतो. हे एक साधे उदाहरण:

\begin{prop}
$x \in B$ असे समजा. मग असा~$A$ आहे की $A \subseteq B$ आणि $A \neq \emptyset$.
\end{prop}

\begin{proof}
$x \in B$ असे गृहीत धरा. $A = \{x\}$ असू द्या.
\begin{quote}
येथे~$A$ संच त्याच्या घटक ची यादी देऊन परिभाषित केला आहे. $x$ ही वस्तू आहे असे गृहीत धरले आहे; आणि सांत यादीत घटक देऊन संच तयार करता येतो. म्हणून येथे~$A$ संच परिभाषित करण्यात यश आले आहे हे वेगळे दाखवण्याची गरज नाही. मात्र $A$ मध्ये विधानाने अपेक्षित केलेले गुणधर्म आहेत हे अद्याप दाखवावे लागते. त्याशिवाय सिद्धता पूर्ण होत नाही!{}
\end{quote}
$x \in A$ असल्यामुळे $A \neq \emptyset$.
\begin{quote}
याचा आधार म्हणजे $A$ ची $\{x\}$ अशी व्याख्या आणि $x \in \{x\}$ तसेच $x \notin \emptyset$ या स्पष्ट गोष्टी.
\end{quote}
$x$ हा~$\{x\}$ मधील एकमेव घटक आहे आणि $x \in B$ आहे. म्हणून~$A$ मधील प्रत्येक घटक हा~$B$ मधील घटक देखील असतो. ~$ \subseteq$ च्या व्याख्येनुसार $A \subseteq B$.
\end{proof}

\subsection{अस्तित्ववाची विधानांचा वापर}

एखादे अस्तित्ववाची विधान सत्य आहे हे माहीत आहे असे समजा: ते सिद्ध केलेले असेल किंवा वापरता येण्याजोगे गृहीतक असेल. उदाहरणार्थ, “काही~$x$ साठी $x \in A$” किंवा “असा $x \in A$ आहे”. सिद्धतेत हे वापरायचे असेल, तर गृहीतकानुसार अस्तित्वात असलेल्या वस्तूंपैकी एका वस्तूला नाव दिले आहे असे धरता येते. $A$ मध्ये किमान एक वस्तू असल्यामुळे त्या नावाने दर्शवता येण्याजोग्या वस्तू आहेत. एखादी वस्तू नेमकी निवडून दाखवता येईलच असे नाही किंवा ती $\in A$ आहे यापलीकडे तिचे वर्णन करता येईलच असे नाही. पण सिद्धतेपुरते एका अशा वस्तूला नाव देता येते. आधी वापरलेले किंवा गृहीतकांत आलेले नाव पुन्हा वापरू नये; तसे केल्यास चुका होऊ शकतात. सिद्धतेत “काही $x$ साठी $x \in A$” यावरून “$a \in A$ असू द्या” असे लिहून हे दर्शवतो. आता~$a$ विषयी विचार करता येतो आणि इतर अनुमत गृहीतके वापरून $p$ हा निष्कर्ष मिळवता येतो. $p$ मध्ये~$a$ चा उल्लेख उरलेला नसेल आणि या नव्या नावाशी संबंधित अतिरिक्त, मुक्त न केलेले गृहीतक वापरलेले नसेल, तर $p$ मिळवण्यासाठीचे तात्पुरते $a \in A$ हे गृहीतक मुक्त करता येते. त्यामुळे “काही $x$ साठी $x \in A$” आणि वापरलेली इतर अनुमत गृहीतके यांवरून निष्कर्ष मिळतो.

\begin{prop}
$A \neq \emptyset$ असेल, तर $A \cup B \neq \emptyset$.
\end{prop}

\begin{proof}
$A \neq \emptyset$ असे समजा. म्हणून काही $x$ साठी $x \in A$.
\begin{quote}
येथे प्रथम विधानाचे गृहीतक फक्त पुन्हा मांडले आहे. या गृहीतकात, म्हणजे $A \neq \emptyset$ यात, एक अस्तित्ववाची विधान दडलेले आहे. काही व्याख्या उलगडल्यावर ते दिसते. $=$ च्या व्याख्येनुसार $A = \emptyset$ हे तेव्हाच आणि फक्त तेव्हा सत्य असते जेव्हा प्रत्येक $x \in A$ हा $\in \emptyset$ देखील असतो आणि प्रत्येक $x \in \emptyset$ हा $\in A$ देखील असतो. दोन्ही बाजूंचे नकरण केल्यावर हे मिळते: $A \neq \emptyset$ तेव्हाच आणि फक्त तेव्हा जेव्हा काही $x \in A$ हा $\notin \emptyset$ असतो किंवा काही $x \in \emptyset$ हा $\notin A$ असतो. कोणतीही वस्तू $\in \emptyset$ नसल्यामुळे दुसरा विकल्पघटक कधीही सत्य असू शकत नाही. तसेच “$x \in A$ आणि $x \notin \emptyset$” याचा अर्थ फक्त $x \in A$ इतकाच उरतो. म्हणून $A \neq \emptyset$ हे काही $x$ साठी $x \in A$ असते तेव्हाच आणि फक्त तेव्हाच सत्य असते. हे अस्तित्ववाची विधान आहे. आता~$A$ मधील घटक पैकी एकाला नाव देऊन या अस्तित्ववाची विधानाचा वापर करू:
\end{quote}
$a \in A$ असू द्या.
\begin{quote}
आता~$\in A$ असलेल्या एका वस्तूला नाव दिले आहे. पुढे~$a$ विषयी युक्तिवाद करू; पण केवळ $a \in A$ आहे एवढेच गृहीत धरू, आणखी काही नाही, याची काळजी घेऊ:
\end{quote}
$a \in A$ असल्यामुळे~$\cup$ च्या व्याख्येवरून $a \in A \cup B$. म्हणून काही $x$ साठी $x \in A \cup B$; म्हणजे $A \cup B \neq \emptyset$.
\begin{quote}
शेवटच्या पायरीत “$a \in A \cup B$” यावरून “काही $x$ साठी $x \in A \cup B$” असा निष्कर्ष काढला. त्यात आता $a$ चा उल्लेख नाही. म्हणून “काही $x$ साठी $x \in A \cup B$” हे केवळ “काही $x$ साठी $x \in A$” यावरून मिळते हे समजते. पण याचाच अर्थ $A \cup B \neq \emptyset$ असा होतो.
\end{quote}
\end{proof}

आपण केले तसे “$x$” सारखे बद्ध चल आणि $a$ सारखी तात्पुरती नावे वेगळी ठेवणे हा चांगला सराव आहे. प्रत्यक्षात मात्र अनेकदा असे करत नाही आणि $x$ हेच वापरतो:
\begin{quote}
$A \neq \emptyset$ असे समजा; म्हणजे असा $x \in A$ आहे. $\cup$ च्या व्याख्येनुसार $x \in A \cup B$. म्हणून $A \cup B \neq \emptyset$.
\end{quote}
मात्र असे करताना वेगवेगळ्या अस्तित्ववाची विधानांसाठी वेगवेगळी $x$ आणि $y$ नावे वापरतो आहोत याची विशेष काळजी घ्यावी. उदाहरणार्थ, पुढील युक्तिवाद ही “$A \neq \emptyset$ आणि $B \neq \emptyset$ असतील, तर $A \cap B \neq \emptyset$” याची योग्य सिद्धता \emph{नाही}; ते विधानच खरे नाही.
\begin{quote}
$A \neq \emptyset$ आणि $B \neq \emptyset$ असे समजा. म्हणून काही $x$ साठी $x \in A$ आणि काही $x$ साठी $x \in B$ देखील असते. $x \in A$ आणि $x \in B$ असल्यामुळे~$\cap$ च्या व्याख्येनुसार $x \in A \cap B$. म्हणून $A \cap B \neq \emptyset$.
\end{quote}
चुकीची पायरी कुठे येते ते ओळखता येईल का? हा निष्कर्ष का खरा नाही हे स्पष्ट करता येईल का?












\section{एक उदाहरण}\label{mth:prf:ex1:sec}

पहिले उदाहरण म्हणजे संचांच्या संयोग आणि छेदाविषयीचे पुढील साधे तथ्य. त्यातून व्याख्या उलगडणे, संधियोग आणि वैश्विक विधाने सिद्ध करणे, तसेच प्रकरणांनुसार सिद्धता करणे या पद्धती स्पष्ट होतील.

\begin{prop}
कोणत्याही $A$, $B$ आणि $C$ संचांसाठी, $A \cup (B \cap C) = (A \cup B) \cap (A \cup C)$.
\end{prop}

हे सिद्ध करू!{}

\begin{proof}
कोणत्याही $A$, $B$ आणि $C$ संचांसाठी $A \cup (B \cap C) = (A \cup B) \cap (A \cup C)$ दाखवायचे आहे.
\begin{quote}
प्रथम विधानातील “$=$” ची व्याख्या उलगडू. संच एकरूप आहेत हे सिद्ध करणे म्हणजे त्यांचे घटक तेच आहेत हे दाखवणे, हे आठवा. म्हणजे $A \cup (B \cap C)$ मधील सर्व घटक हे $(A \cup B) \cap (A \cup C)$ मधील घटक देखील असले पाहिजेत आणि उलटही तसेच असले पाहिजे. “उलटही तसेच” म्हणजे $(A \cup B) \cap (A \cup C)$ मधील प्रत्येक घटक हा $A \cup (B \cap C)$ मधील घटक देखील असला पाहिजे. म्हणून व्याख्या उलगडल्यावर संधियोग सिद्ध करावा लागतो असे दिसते. हे नोंदवू:
\end{quote}
व्याख्येनुसार, $A \cup (B \cap C) = (A \cup B) \cap (A \cup C)$ हे तेव्हाच आणि फक्त तेव्हा सत्य असते जेव्हा $A \cup (B \cap C)$ मधील प्रत्येक घटक हा $(A \cup B) \cap (A \cup C)$ मधील घटक देखील असतो आणि $(A \cup B) \cap (A \cup C)$ मधील प्रत्येक घटक हा $A \cup (B \cap C)$ मधील घटक असतो.
\begin{quote}
हा संधियोग असल्यामुळे प्रत्येक संधिघटक वेगळा सिद्ध करावा लागेल. पहिल्यापासून सुरुवात करू: $A \cup (B \cap C)$ मधील प्रत्येक घटक हा $(A \cup B) \cap (A \cup C)$ मधील घटक देखील असतो हे सिद्ध करू.

हे वैश्विक विधान आहे. म्हणून $A \cup (B \cap C)$ मधील एक स्वेच्छ घटक घेऊन तो $(A \cup B) \cap (A \cup C)$ मधील घटक देखील असला पाहिजे हे दाखवू. या स्वेच्छ घटक साठी एखादे चल निवडू; उदाहरणार्थ,~$z$. सिद्धता पुढे अशी चालते:
\end{quote}
प्रथम $A \cup (B \cap C)$ मधील प्रत्येक घटक हा $(A \cup B) \cap (A \cup C)$ मधील घटक देखील असतो हे सिद्ध करू. $z \in A \cup (B \cap C)$ असू द्या. $z \in (A \cup B) \cap (A \cup C)$ दाखवायचे आहे.
\begin{quote}
आता $\cup$ आणि~$\cap$ यांच्या व्याख्या उलगडायच्या आहेत. उदाहरणार्थ, $\cup$ ची व्याख्या अशी आहे: $A \cup B = \Setabs{z}{z \in A \text{ किंवा } z \in B}$. “$A \cup (B \cap C)$” ला ही व्याख्या लावताना व्याख्येतील “$B$” ची भूमिका आता “$B \cap C$” बजावते. म्हणून $A \cup (B \cap C) = \Setabs{z}{z \in A \text{ किंवा } z \in B \cap C}$. त्यामुळे $z \in A \cup (B \cap C)$ या गृहीतकाचा अर्थ $z \in \Setabs{z}{z \in A \text{ किंवा } z \in B \cap C}$ असा होतो. तसेच $z \in \Setabs{z}{\dots z\dots}$ हे \dots $z$ \dots{} असते तेव्हाच आणि फक्त तेव्हाच सत्य असते. म्हणजे येथे $z \in A$ किंवा $z \in B \cap C$ असते.
\end{quote}
$\cup$ च्या व्याख्येनुसार $z \in A$ किंवा $z \in B \cap C$ असते.
\begin{quote}
हा विकल्पयोग असल्यामुळे प्रकरणांनुसार सिद्धता उपयुक्त ठरेल. दोन्ही प्रकरणे घेऊन प्रत्येक प्रकरणात अपेक्षित निष्कर्ष---म्हणजे “$z \in (A \cup B) \cap (A \cup C)$”---मिळतो हे दाखवू.
\end{quote}
प्रकरण 1: $z \in A$ असे समजा.
\begin{quote}
या गृहीतकांवरून आणखी फार काही मिळत नाही. म्हणून अपेक्षित निष्कर्षात कोणती रचना आहे ते पाहू. $z \in (A \cup B) \cap (A \cup C)$ दाखवायचे आहे. $\cap$ च्या व्याख्येनुसार $z \in (A \cup B) \cap (A \cup C)$ दाखवण्यासाठी तो $(A \cup B)$ आणि $(A \cup C)$ या दोन्हींत आहे हे दाखवावे लागते. पण $z \in A \cup B$ हे $z \in A$ किंवा $z \in B$ असते तेव्हाच आणि फक्त तेव्हाच सत्य असते. तसेच प्रकरण~1 चे गृहीतक $z \in A$ आधीच उपलब्ध आहे. हाच युक्तिवाद $B$ ऐवजी $C$ वापरून केल्यास $z \in A \cup C$ मिळते. हा विचार अपेक्षित निष्कर्षापासून मागे जात केलेला होता; आता सिद्धतेत आवश्यक असलेल्या पुढच्या दिशेने तो नोंदवू.
\end{quote}
$z \in A$ असल्यामुळे $z \in A$ किंवा $z \in B$. म्हणून~$\cup$ च्या व्याख्येनुसार $z \in A \cup B$. त्याचप्रमाणे $z \in A \cup C$. पण~$\cap$ च्या व्याख्येनुसार याचाच अर्थ $z \in (A \cup B) \cap (A \cup C)$ असा होतो.
\begin{quote}
यामुळे प्रकरणांनुसार सिद्धतेचे पहिले प्रकरण पूर्ण झाले. आता $z \in B \cap C$ असलेल्या दुसऱ्या प्रकरणात निष्कर्ष मिळवायचा आहे.
\end{quote}
प्रकरण 2: $z \in B \cap C$ असे समजा.
\begin{quote}
पुन्हा दोन संचांचा छेद हाताळत आहोत. ~$ \cap$ ची व्याख्या वापरू:
\end{quote}
$z \in B \cap C$ असल्यामुळे~$\cap$ च्या व्याख्येनुसार $z$ हा $B$ आणि $C$ या दोन्हींचा घटक असला पाहिजे.
\begin{quote}
पुन्हा अपेक्षित निष्कर्ष पाहू. $z$ हा $(A \cup B)$ आणि $(A \cup C)$ या दोन्हींत आहे हे दाखवायचे आहे. पुन्हा उत्तर लगेच मिळते.
\end{quote}
$z \in B$ असल्यामुळे $z \in (A \cup B)$. $z \in C$ असल्यामुळे $z \in (A \cup C)$ देखील. म्हणून $z \in (A \cup B) \cap (A \cup C)$.
\begin{quote}
येथे पुन्हा $\cup$ आणि $\cap$ यांच्या व्याख्या वापरल्या. पण त्या व्याख्या आधीच आठवून घेतल्या आहेत आणि $z$ हा दोन संचांपैकी एका संचात असेल तर त्यांच्या संयोगात असतो, हेही आधीच दाखवले आहे. म्हणून आता प्रत्येक वापराचा तितका तपशील लिहिण्याची गरज नाही.

प्रकरणांनुसार सिद्धतेचे दुसरे प्रकरण पूर्ण झाले. म्हणून पहिला अपेक्षित निष्कर्ष आता सांगता येतो.
\end{quote}
म्हणून $z \in A \cup (B \cap C)$ असेल, तर $z \in (A \cup B) \cap (A \cup C)$.
\begin{quote}
आता उलटी दिशा दाखवायची आहे: $(A \cup B) \cap (A \cup C)$ मधील प्रत्येक घटक हा $A \cup (B \cap C)$ मधील घटक आहे. आधीप्रमाणे, पहिल्या संचातील एक स्वेच्छ घटक घेतला आहे असे गृहीत धरून तो दुसऱ्या संचात असला पाहिजे हे दाखवून हे वैश्विक विधान सिद्ध करू. काय करणार आहोत ते मांडू.
\end{quote}
आता $z \in (A \cup B) \cap (A \cup C)$ असे गृहीत धरा. $z \in A \cup (B \cap C)$ दाखवायचे आहे.
\begin{quote}
आता $z \in (A \cup B) \cap (A \cup C)$ या गृहीतकावरून विचार करत आहोत. सिद्धतेच्या पहिल्या भागात वापरलेले~$z$ हेच येथे वापरल्यामुळे फार गोंधळ होत नाही अशी आशा आहे. तो भाग पूर्ण झाल्यावर त्या भागातील तात्पुरती गृहीतके लागू राहिली नाहीत; म्हणून आता $z$ विषयी नवी गृहीतके घेता येतात. हे गोंधळात टाकत असेल, तर पुढील भागात $z$ ऐवजी वेगळे चल वापरा.

~$\cap$ च्या व्याख्येनुसार $z$ हा $A \cup B$ आणि $A \cup C$ या दोन्हींत आहे. पुढे $\cup$ ची व्याख्या उलगडल्यावर असे मिळते: $z \in A$ किंवा $z \in B$, आणि $z \in A$ किंवा $z \in C$ देखील. पुन्हा प्रकरणांनुसार सिद्धता दिसते; पण “आणि” मुळे गोंधळ होऊ शकतो. दोन्ही अटींऐवजी एकच अट पुरेशी आहे---$z$ हा $A$, $B$ किंवा $C$ पैकी एखाद्यात आहे---असे वाटू शकते. पण ती अदलाबदल चुकीची ठरेल. म्हणून काळजीपूर्वक एकेक विकल्पयोग पाहू.
\end{quote}
$\cap$ च्या व्याख्येनुसार $z \in A \cup B$ आणि $z \in A \cup C$. $\cup$ च्या व्याख्येनुसार $z \in A$ किंवा $z \in B$. प्रकरणे वेगळी पाहू.
\begin{quote}
पहिल्या विकल्पयोगावर लक्ष केंद्रित केल्यामुळे $A \cup C$ उलगडून मिळणारा दुसरा विकल्पयोग अजून नोंदवलेला नाही. आत्ता त्याची गरजही नाही. पहिले प्रकरण $z \in A$ हे आहे. संचाचा घटक हा त्या संचाचा दुसऱ्या कोणत्याही संचाशी संयोग केल्यावर मिळणाऱ्या संचाचा घटक देखील असतो. म्हणून प्रकरण~1 सोपे आहे:
\end{quote}
प्रकरण 1: $z \in A$ असे समजा. यावरून $z \in A \cup (B \cap C)$ मिळते.
\begin{quote}
आता दुसरे प्रकरण, $z \in B$, पाहू. येथे दुसऱ्या $\cup$ ची व्याख्या उलगडून आणखी एक प्रकरणांनुसार सिद्धता करू:
\end{quote}
प्रकरण 2: $z \in B$ असे समजा. $z \in A \cup C$ असल्यामुळे $z \in A$ किंवा $z \in C$. आणखी प्रकरणे वेगळी पाहू:

प्रकरण 2a: $z \in A$. मग पुन्हा $z \in A \cup (B \cap C)$.
\begin{quote}
हे थोडे विचित्र वाटेल: या प्रकरणात~$z \in B$ हे गृहीतक वापरण्याची गरजच पडली नाही. पण त्यात काही अडचण नाही.
\end{quote}
प्रकरण 2b: $z \in C$. मग $z \in B$ आणि $z \in C$. म्हणून $z \in B \cap C$ आणि परिणामी $z \in A \cup (B \cap C)$.
\begin{quote}
दोन्ही प्रकरणांनुसार सिद्धता पूर्ण झाल्या; त्यामुळे दुसरा अर्धा भागही पूर्ण झाला.
\end{quote}
म्हणून $z \in (A \cup B) \cap (A \cup C)$ असेल, तर $z \in A \cup (B \cap C)$.
\end{proof}












\section{आणखी एक उदाहरण}\label{mth:prf:ex2:sec}

\begin{prop}
$A \subseteq C$ असेल, तर $A \cup (C \setminus A) = C$.
\end{prop}

\begin{proof}
$A \subseteq C$ असे समजा. $A \cup (C \setminus A) = C$ दाखवायचे आहे.
\begin{quote}
सुरुवातीला हे सशर्त विधान आहे हे लक्षात घेऊ. त्याचे वैश्विक संख्यापन अप्रकट आहे: हे विधान सर्व $A$ आणि $C$ संचांसाठी खरे आहे. म्हणून $A$ आणि $C$ हे स्वेच्छ संचांचे चल आहेत. असे विधान सिद्ध करण्यासाठी पूर्वांग गृहीत धरून उत्तरांग सिद्ध करतो.

आता $A \subseteq C$ हे गृहीतक वापरू. ~$\subseteq$ ची व्याख्या उलगडू: गृहीतकाचा अर्थ असा की~$A$ मधील सर्व घटक हे~$C$ मधील घटक देखील आहेत. हे नोंदवू; संपूर्ण सिद्धतेत वापरणार असलेले हे महत्त्वाचे तथ्य आहे.
\end{quote}
~$\subseteq$ च्या व्याख्येनुसार, $A \subseteq C$ असल्यामुळे प्रत्येक $z$ साठी $z \in A$ असेल, तर $z \in C$ असते.
\begin{quote}
गृहीतकातील सर्व व्याख्या उलगडल्या आहेत. आता निष्कर्षाकडे वळू. $A \cup (C \setminus A) = C$ दाखवायचे आहे. म्हणून मागच्या उदाहरणाप्रमाणे सिद्धतेची मांडणी करू: $A \cup (C \setminus A)$ मधील प्रत्येक घटक हा~$C$ मधील घटक आहे आणि उलट $C$ मधील प्रत्येक घटक हा $A \cup (C \setminus A)$ मधील घटक आहे, हे दाखवू. हे थोडक्यात असे लिहिता येईल: $A \cup (C \setminus A) \subseteq C$ आणि $C \subseteq A \cup (C \setminus A)$. येथे व्याख्या उलगडण्याच्या उलट क्रिया करत आहोत; पण त्यामुळे सिद्धता वाचणे थोडे सोपे होते. हा संधियोग असल्यामुळे दोन्ही भाग सिद्ध करावे लागतात. पहिला भाग---म्हणजे $A \cup (C \setminus A)$ मधील प्रत्येक घटक हा~$C$ मधील घटक आहे---दाखवण्यासाठी, स्वेच्छ~$z$ साठी $z \in A \cup (C \setminus A)$ गृहीत धरून $z \in C$ दाखवू. $\cup$ च्या व्याख्येनुसार $z \in A \cup (C \setminus A)$ यावरून $z \in A$ किंवा $z \in C \setminus A$ असा निष्कर्ष काढता येतो. आता तुम्हाला या पद्धतीचा अंदाज येत असेल.
\end{quote}
$A \cup (C \setminus A) = C$ हे तेव्हाच आणि फक्त तेव्हा सत्य असते जेव्हा $A \cup (C \setminus A) \subseteq C$ आणि $C \subseteq (A \cup (C \setminus A))$ असतात. प्रथम $A \cup (C \setminus A) \subseteq C$ सिद्ध करू. $z \in A \cup (C \setminus A)$ असू द्या. मग $z \in A$ किंवा $z \in (C \setminus A)$.
\begin{quote}
एक विकल्पयोग मिळाला आहे; त्यावरून $z \in C$ सिद्ध करायचे आहे. यासाठी प्रकरणांनुसार सिद्धता करू.
\end{quote}
प्रकरण 1: $z \in A$. प्रत्येक $z$ साठी $z \in A$ असेल तर $z \in C$ असते, हे माहीत असल्यामुळे $z \in C$ मिळते.
\begin{quote}
येथे $A \subseteq C$ या गृहीतकावरून मिळालेले आणि आधी नोंदवलेले तथ्य वापरले. पहिले प्रकरण पूर्ण झाले; आता दुसऱ्या प्रकरणाकडे, $z \in (C \setminus A)$, वळू. $C \setminus A$ हे दोन संचांचा \emph{फरक} दर्शवते, हे आठवा:~$C$ मधील जे घटक हे~$A$ मधील घटक नाहीत त्यांचा हा संच आहे. पण $C$ मधील जो घटक हा~$A$ मध्ये नाही तो विशेषतः~$C$ मधील घटक असतोच.
\end{quote}
प्रकरण 2: $z \in (C \setminus A)$. याचा अर्थ $z \in C$ आणि $z \notin A$. म्हणून विशेषतः $z \in C$.
\begin{quote}
पहिली दिशा सिद्ध झाली. आता दुसरी दिशा पाहू. येथे $C \subseteq A \cup (C \setminus A)$ सिद्ध करायचे आहे. म्हणून $z \in C$ गृहीत धरून $z \in A \cup (C \setminus A)$ सिद्ध करू.
\end{quote}
आता $z \in C$ असू द्या. $z \in A$ किंवा $z \in C \setminus A$ दाखवायचे आहे.
\begin{quote}
$A$ मधील सर्व घटक हे $C$ मधील घटक देखील आहेत; आणि $C \setminus A$ हा $C$ मधील, पण $A$ मधील नसलेल्या, घटक चा संच आहे. म्हणून $z$ हा $A$ मध्ये किंवा $C \setminus A$ मध्ये असतो असे दिसते. निष्कर्ष का खरा आहे हे आधीच माहीत नसेल, तर हा युक्तिवाद थोडा अस्पष्ट वाटेल. तो पायरीपायरीने सिद्ध करणे अधिक चांगले. सिद्धता न देता मांडता येणारे एक साधे तथ्य वापरल्यास मदत होईल: $z \in A$ किंवा $z \notin A$. याला “विमध्य सिद्धांत” म्हणतात: कोणत्याही~$p$ विधानासाठी $p$ सत्य असते किंवा त्याचे नकरण सत्य असते. येथे $p$ हे $z \in A$ हे विधान आहे. हा विकल्पयोग असल्यामुळे पुन्हा प्रकरणांनुसार सिद्धता वापरता येते.
\end{quote}
$z \in A$ किंवा $z \notin A$. पहिल्या प्रकरणात $z \in A \cup (C \setminus A)$. दुसऱ्या प्रकरणात $z \in C$ आणि $z \notin A$. म्हणून $z \in C \setminus A$. पण मग $z \in A \cup (C \setminus A)$.
\begin{quote}
सिद्धता पूर्ण झाली: $A \cup (C \setminus A) = C$ दाखवले आहे.
\end{quote}
\end{proof}












\section{व्याघाताने सिद्धता}\label{mth:prf:con:sec}

व्याघाताने सिद्धता हा प्रथमतः नकारात्मक विधाने सिद्ध करण्यासाठी वापरला जाणारा अनुमानाचा नमुना आहे. एखादे~$p$ विधान \emph{असत्य} आहे हे दाखवायचे आहे असे समजा; म्हणजे~$\lnot p$ दाखवायचे आहे. यासाठी सर्वात उपयोगी मार्ग असा: (a) $p$~सत्य आहे असे समजा आणि (b) या गृहीतकावरून असत्य असल्याचे माहीत असलेली एखादी गोष्ट मिळते हे दाखवा. “असत्य असल्याचे माहीत असलेली गोष्ट” ही $p$ शीच व्याघात करणारा निष्कर्ष असेल किंवा संपूर्ण विधानाच्या दुसऱ्या एखाद्या गृहीतकाशी व्याघात करणारा निष्कर्ष असेल. उदाहरणार्थ, “$q$ असेल, तर $\lnot p$” सिद्ध करताना $q$~सत्य आहे असे गृहीत धरून त्यावरून~$\lnot p$ सिद्ध करतो. $\lnot p$ व्याघाताने सिद्ध करायचे असेल, तर~$q$ बरोबर $p$ हेदेखील गृहीत धरायचे. $p$ वरून $\lnot q$ सिद्ध करता आले, तर~$p$ या गृहीतकावरून दुसऱ्या~$q$ गृहीतकाशी व्याघात करणारा निष्कर्ष मिळतो हे दाखवले आहे. कारण $q$~आणि $\lnot q$ दोन्ही सत्य असू शकत नाहीत. अर्थात, व्याघात सिद्ध करताना इतर अनुमानाचे नमुने वापरावे लागतील आणि व्याख्याही उलगडाव्या लागतील. एक उदाहरण पाहू.

\begin{prop}
$A \subseteq B$ आणि $B = \emptyset$ असतील, तर $A$ मध्ये कोणतेही घटक नसतात.
\end{prop}

\begin{proof}
$A \subseteq B$ आणि $B = \emptyset$ असे समजा. $A$ मध्ये कोणतेही घटक नाहीत हे दाखवायचे आहे.
\begin{quote}
हे सशर्त विधान असल्यामुळे पूर्वांग गृहीत धरून उत्तरांग सिद्ध करायचे आहे. उत्तरांग असे आहे: $A$ मध्ये कोणतेही घटक नाहीत. हे आणखी स्पष्टपणे असे मांडता येते: असा~$x \in A$ आहे हे खरे नाही.
\end{quote}
$A$ मध्ये कोणतेही घटक नसतात तेव्हाच आणि फक्त तेव्हा जेव्हा असा~$x$ आहे की $x \in A$, हे खरे नसते.
\begin{quote}
म्हणून प्रत्यक्षात~$\lnot p$ हे नकारात्मक विधान सिद्ध करायचे आहे असे ठरले: असा $x \in A$ आहे हे खरे नाही. व्याघाताने सिद्धता करण्यासाठी त्याच्याशी संबंधित~$p$ हे होकारात्मक विधान---म्हणजे असा $x \in A$ आहे---गृहीत धरून त्यावरून व्याघात सिद्ध करावा लागेल. “व्याघात मिळवण्यासाठी असे गृहीत धरा” किंवा नुसते “उलट समजा” असे लिहून~$p$ हे गृहीतक मांडले, की आपण व्याघाताने सिद्धता करत आहोत हे दर्शवता येते.
\end{quote}
उलट समजा: असा $x \in A$ आहे.
\begin{quote}
व्याघात मिळवण्यासाठी हे नवे गृहीतक वापरू. आणखी दोन गृहीतके आहेत: $A \subseteq B$ आणि $B = \emptyset$. पहिल्यावरून $x \in B$ मिळते:
\end{quote}
$A \subseteq B$ असल्यामुळे $x \in B$.
\begin{quote}
पण $B = \emptyset$ असल्यामुळे $B$ मधील प्रत्येक घटक---उदाहरणार्थ $x$---हा~$\emptyset$ मधील घटक देखील असला पाहिजे.
\end{quote}
$B = \emptyset$ असल्यामुळे $x \in \emptyset$. हा व्याघात आहे; कारण व्याख्येनुसार $\emptyset$ मध्ये कोणतेही घटक नसतात.
\begin{quote}
यानेच सिद्धता पूर्ण होते. घेतलेल्या गृहीतकांवरून आवश्यक व्याघात मिळाला आहे; म्हणजे ती सर्व गृहीतके एकाच वेळी खरी असू शकत नाहीत. पहिली दोन गृहीतके ($A \subseteq B$ आणि $B = \emptyset$) कायम ठेवलेली असल्यामुळे शेवटी आणलेले गृहीतक---असा $x \in A$ आहे---असत्य असले पाहिजे. मात्र पूर्ण तपशील द्यायचा असल्यास हे स्पष्ट लिहिता येते.
\end{quote}
म्हणून असा $x \in A$ आहे हे गृहीतक असत्य असले पाहिजे. त्यामुळे व्याघाताने सिद्धतेवरून $A$ मध्ये कोणतेही घटक नाहीत.
\end{proof}

शास्त्रीय तर्कशास्त्रात प्रत्येक होकारात्मक विधान एका नकारात्मक विधानाशी सममूल्य असते: $p$ हे $\lnot\lnot p$ असते तेव्हाच आणि फक्त तेव्हाच सत्य असते. म्हणून व्याघाताने सिद्धता वापरून होकारात्मक विधानेही “अप्रत्यक्षपणे” सिद्ध करता येतात. ~$p$ सिद्ध करण्यासाठी त्याला $\lnot\lnot p$ हे नकारात्मक विधान समजा. $\lnot p$ वरून व्याघात सिद्ध करता आला, तर व्याघाताने सिद्धतेद्वारे $\lnot\lnot p$ प्रस्थापित होते आणि म्हणून~$p$ मिळते.

मागील उदाहरणात नकारात्मक विधान सिद्ध करायचे होते: $A$ मध्ये कोणतेही घटक नाहीत. म्हणून व्याघाताने सिद्धतेसाठी घेतलेले गृहीतक---म्हणजे असा $x \in A$ आहे---हे होकारात्मक होते. त्यामुळे विचार करण्यासाठी तात्पुरता $x \in A$ मिळाला; त्याच्याविषयी विचार पुढे नेत $x \in \emptyset$ पर्यंत पोहोचलो.

होकारात्मक विधान अप्रत्यक्षपणे सिद्ध करताना व्याघात मिळवण्यासाठी नकारात्मक गृहीतक घ्यावे लागते. पण अनेकदा होकारात्मक विधान नकारात्मक रूपात आणि नकारात्मक विधान होकारात्मक रूपात सहज पुन्हा मांडता येते. मागच्या सिद्धतेत “$A$ मध्ये कोणतेही घटक नाहीत” या नकारात्मक उत्तरांगाऐवजी “$A = \emptyset$” सिद्ध केले असते, तरी सिद्धता मूलतः तशीच असती. $=$ च्या व्याख्येनुसार “$A = \emptyset$” हे सामान्य विधान आहे; कारण ते उलगडल्यावर “~$A$ मधील प्रत्येक घटक हा~$\emptyset$ मधील घटक आहे आणि उलटही तसेच आहे” असे मिळते. पण ते “असा $x \in A$ आहे हे खरे नाही” या नकारात्मक विधानाशी सममूल्य आहे हे सहज दिसते.

म्हणून कधी~$p$ प्रत्यक्ष सिद्ध करण्यापेक्षा $\lnot p$ हे गृहीतक घेऊन काम करणे सोपे असते. पुढील उदाहरणांप्रमाणे प्रत्यक्ष सिद्धता तितकीच सोपी किंवा अधिक सोपी असली, तरी काही जण अप्रत्यक्ष मार्ग पसंत करतात. द्विनकरणामुळे गोंधळ होत असेल, तर व्याघाताने सिद्धता म्हणजे \emph{विरुद्ध} विधान गृहीत धरून व्याघात सिद्ध करणे, असा विचार करा. म्हणून $\lnot p$ व्याघाताने सिद्ध करणे म्हणजे~$p$ गृहीत धरून व्याघात सिद्ध करणे; आणि~$p$ व्याघाताने सिद्ध करणे म्हणजे~$\lnot p$ वरून व्याघात सिद्ध करणे.

\begin{prop}
$A \subseteq A \cup B$.
\end{prop}

\begin{proof}
$A \subseteq A \cup B$ दाखवायचे आहे.
\begin{quote}
प्रथमदर्शनी हे होकारात्मक विधान आहे: प्रत्येक $x \in A$ हा $A \cup B$ मध्येही असतो. त्याचे नकरण असे: काही $x \in A$ हा $\notin A \cup B$ असतो. म्हणून हे नकरण गृहीत धरून त्यावरून व्याघात मिळतो असे दाखवल्यास विधान अप्रत्यक्षपणे सिद्ध होते.
\end{quote}
उलट समजा; म्हणजे $A \nsubseteq A \cup B$.
\begin{quote}
$A \subseteq A \cup B$ ची व्याख्या अशी आहे: प्रत्येक $x \in A$ हा $\in A \cup B$ देखील असतो. $A \nsubseteq A \cup B$ चा अर्थ समजण्यासाठी उलगडलेल्या व्याख्येवर काही मूलभूत तार्किक क्रिया कराव्या लागतात: प्रत्येक $x \in A$ हा $\in A \cup B$ देखील असतो हे असत्य असते तेव्हाच आणि फक्त तेव्हा जेव्हा असा \emph{काही}~$x \in A$ असतो की तो $\notin A \cup B$ असतो. येथे फार काळजी घ्यावी: “सर्व $A$ हे $B$ आहेत” अशा विधानाचे नकरण चुकीचे उलगडल्यामुळे अनेक विद्यार्थ्यांच्या व्याघाताने सिद्धता करण्याच्या प्रयत्नांत चूक होते. म्हणजे $A \nsubseteq A \cup B$ हे तेव्हाच आणि फक्त तेव्हा सत्य असते जेव्हा असा $x$ आहे की $x \in A$ आणि $x \notin A \cup B$. यापुढे सिद्धता सोपी आहे.
\end{quote}
म्हणून असा $x \in A$ आहे की $x \notin A \cup B$. $\cup$ च्या व्याख्येनुसार $x \in A \cup B$ हे $x \in A$ किंवा $x \in B$ असते तेव्हाच आणि फक्त तेव्हाच सत्य असते. $x \in A$ असल्यामुळे $x \in A \cup B$ मिळते. हा $x \notin A \cup B$ या गृहीतकाशी व्याघात आहे.
\end{proof}

\begin{prob}
$A \cap B \subseteq A$ हे \emph{अप्रत्यक्षपणे} सिद्ध करा.
\end{prob}

\begin{prop}
$A \subseteq B$ आणि $B \subseteq C$ असतील, तर $A \subseteq C$.
\end{prop}

\begin{proof}
$A \subseteq B$ आणि $B \subseteq C$ असे समजा. $A \subseteq C$ दाखवायचे आहे.
\begin{quote}
अप्रत्यक्ष मार्ग वापरू: जे प्रस्थापित करायचे आहे त्याचे नकरण गृहीत धरू.
\end{quote}
उलट समजा; म्हणजे $A \nsubseteq C$.
\begin{quote}
आधीप्रमाणे, $A \nsubseteq C$ हे तेव्हाच आणि फक्त तेव्हा सत्य असते जेव्हा प्रत्येक $x \in A$ हा $\in C$ देखील असतो हे खरे नसते; म्हणजे काही $x \in A$ हा $\notin C$ असतो. सराव झाल्यावर अशी नकरणे उलगडण्यासाठी फार विचार करावा लागणार नाही.
\end{quote}
म्हणजे असा~$x$ आहे की $x \in A$ आणि $x \notin C$.
\begin{quote}
आता यावरून व्याघात मिळवता येईल. अर्थात, त्यासाठी इतर दोन गृहीतके वापरावी लागतील.
\end{quote}
$A \subseteq B$ असल्यामुळे $x \in B$. $B \subseteq C$ असल्यामुळे $x \in C$. पण हा $x \notin C$ शी व्याघात आहे.
\end{proof}

\begin{prop}
$A \cup B = A \cap B$ असेल, तर $A = B$.
\end{prop}

\begin{proof}
$A \cup B = A \cap B$ असे समजा. $A = B$ दाखवायचे आहे.
\begin{quote}
आता सुरुवातीची पद्धत परिचित आहे:
\end{quote}
व्याघात मिळवण्यासाठी $A \neq B$ असे गृहीत धरा.
\begin{quote}
व्याघाताने सिद्धतेसाठीचे गृहीतक $A \neq B$ आहे. $A = B$ हे $A \subseteq B$ आणि $B \subseteq A$ असतात तेव्हाच आणि फक्त तेव्हाच सत्य असल्यामुळे, $A \neq B$ हे $A \nsubseteq B$ \emph{किंवा} $B \nsubseteq A$ असते तेव्हाच आणि फक्त तेव्हाच सत्य असते. नकरणांवरील क्रिया काळजीपूर्वक करणे किती महत्त्वाचे आहे ते पाहा!{} या विकल्पयोगावरून व्याघात सिद्ध करण्यासाठी प्रकरणांनुसार सिद्धता वापरून प्रत्येक प्रकरणात व्याघात मिळतो हे दाखवू.
\end{quote}
$A \neq B$ हे $A \nsubseteq B$ किंवा $B \nsubseteq A$ असते तेव्हाच आणि फक्त तेव्हाच सत्य असते. प्रकरणे वेगळी पाहू.
\begin{quote}
पहिल्या प्रकरणात $A \nsubseteq B$ गृहीत धरतो; म्हणजे काही $x$ साठी $x \in A$, पण $\notin B$. $A \cap B$ ची व्याख्या $A$ आणि $B$ यांच्या समान घटक चा संच अशी आहे. म्हणून एखादी वस्तू त्यांपैकी एका संचात नसेल, तर ती छेदात नाही. $A \cup B$ हा $A$ बरोबर $B$ घेऊन मिळालेला संच आहे. म्हणून कोणत्याही एका संचातील वस्तू संयोगातही असते. यावरून $x \in A \cup B$, पण $x \notin A \cap B$, असे मिळते; म्हणून $A \cap B \neq A \cup B$.
\end{quote}

प्रकरण 1: $A \nsubseteq B$. मग काही $x$ साठी $x \in A$, पण $x \notin B$. $x \notin B$ असल्यामुळे $x \notin A \cap B$. $x \in A$ असल्यामुळे $x \in A \cup B$. म्हणून $A \cap B \neq A \cup B$; हा $A \cap B = A \cup B$ या गृहीतकाशी व्याघात आहे.

प्रकरण 2: $B \nsubseteq A$. मग काही $y$ साठी $y \in B$, पण $y \notin A$. आधीप्रमाणे $y \in A \cup B$, पण $y \notin A \cap B$. म्हणून $A \cap B \neq A \cup B$; हा पुन्हा $A \cap B = A \cup B$ शी व्याघात आहे.
\end{proof}











\section{सिद्धता वाचणे}\label{mth:prf:rea:sec}

पाठ्यपुस्तकांतील आणि लेखांतील सिद्धतांत आपण आत्तापर्यंत उदाहरणांत दिलेला सर्व तपशील क्वचितच असतो. प्रकरणे वेगळी पाहत आहोत, अप्रत्यक्ष सिद्धता देत आहोत किंवा व्याख्या वापरत आहोत, याकडे लेखक अनेकदा स्वतंत्र लक्ष वेधत नाहीत. म्हणून पाठ्यपुस्तकातील सिद्धता वाचताना ती समजण्यासाठी हा तपशील अनेकदा स्वतः भरून काढावा लागतो. सिद्धतेत कराव्या लागणाऱ्या विविध क्रियांचा सराव होण्यासाठीही हे उपयुक्त आहे. एक उदाहरण पाहू.

\begin{prop}[शोषणाचा नियम]
सर्व $A$, $B$ संचांसाठी,
\[
A \cap (A \cup B) = A
\]
\end{prop}

\begin{proof}
$z \in A \cap (A \cup B)$ असेल, तर $z \in A$. म्हणून $A \cap (A \cup B) \subseteq A$. आता $z \in A$ असे समजा. मग $z \in A \cup B$ देखील; म्हणून $z \in A \cap (A \cup B)$ देखील.
\end{proof}

शोषणाच्या नियमाची वरील सिद्धता फार संक्षिप्त आहे. वापरलेल्या व्याख्यांचा उल्लेख नाही; सिद्ध करण्याआधी “हे सिद्ध करायचे आहे” असे म्हटलेले नाही; इत्यादी. ती उलगडू. सिद्ध केलेले विधान हे कोणत्याही $A$ आणि $B$ संचांविषयीचे सामान्य विधान आहे. सिद्धतेत $A$ किंवा $B$ येतात तेव्हा ते स्वेच्छ संचांचे चल आहेत. सिद्धतेत मिळवलेली सामान्य विधाने ही संचांची एकरूपता सिद्ध करण्यासाठी आवश्यक असलेली विधाने आहेत: समीकरणाच्या डाव्या बाजूतील प्रत्येक घटक हा उजव्या बाजूतील घटक असतो आणि उलटही तसेच असते.

\begin{quote}
“$z \in A \cap (A \cup B)$ असेल, तर $z \in A$. म्हणून $A \cap (A \cup B) \subseteq A$.”
\end{quote}

हा एकरूपतेच्या सिद्धतेचा पहिला अर्धा भाग आहे. त्यात असे दाखवले आहे: स्वेच्छ~$z$ हा डाव्या बाजूतील घटक असेल, तर तो उजव्या बाजूतील घटक देखील असतो; म्हणजे $A \cap (A \cup B) \subseteq A$. $z \in A \cap (A \cup B)$ असे गृहीत धरा. $z$ हा दोन संचांच्या छेदाचा घटक असतो तेव्हाच आणि फक्त तेव्हा जेव्हा तो दोन्ही संचांचा घटक असतो. म्हणून $z \in A$ आणि $z \in A \cup B$ देखील, असा निष्कर्ष काढता येतो. विशेषतः $z \in A$ मिळते; हेच दाखवायचे होते. पहिल्या अर्ध्या भागासाठी इतकेच आवश्यक असल्यामुळे उरलेली सिद्धता दुसऱ्या अर्ध्या भागाची असली पाहिजे हे समजते; म्हणजे $A \subseteq A \cap (A \cup B)$ ची सिद्धता.

\begin{quote}
“आता $z \in A$ असे समजा. मग $z \in A \cup B$ देखील; म्हणून $z \in A \cap (A \cup B)$ देखील.”
\end{quote}

सुरुवातीला $z \in A$ असे गृहीत धरतो; कारण कोणत्याही~$z$ साठी $z \in A$ असेल, तर $z \in A \cap (A \cup B)$ असते हे दाखवत आहोत. $z \in A \cap (A \cup B)$ दाखवण्यासाठी “$\cap$” च्या व्याख्येनुसार (i) $z \in A$ आणि (ii) $z \in A \cup B$ दोन्ही दाखवावे लागतात. येथे (i) हे आपले गृहीतकच आहे; म्हणून त्यासाठी आणखी काही सिद्ध करायचे नाही. म्हणूनच सिद्धतेत त्याचा पुन्हा उल्लेख नाही. (ii) साठी हे आठवा: $z$ हा संचांच्या संयोगाचा घटक असतो तेव्हाच आणि फक्त तेव्हा जेव्हा तो त्यांपैकी किमान एका संचाचा घटक असतो. येथे $z \in A$ आहे आणि $A \cup B$ हा $A$ आणि $B$ यांचा संयोग आहे; म्हणून ही अट पूर्ण होते. त्यामुळे $z \in A \cup B$. (i) $z \in A$ आणि (ii) $z \in A \cup B$ दोन्ही दाखवले आहेत. म्हणून “$\cap$” च्या व्याख्येनुसार $z \in A \cap (A \cup B)$. सिद्धतेत या व्याख्यांचा उल्लेख नाही; त्या वाचकाला पूर्णपणे परिचित आहेत असे मानलेले आहे. तसे नसल्यास मागे जाऊन त्या काय आहेत हे आठवून घ्यावे लागेल. मग $z \in A$ वरून $z \in A \cup B$ का मिळते आणि $z \in A$ तसेच $z \in A \cup B$ यांवरून $z \in A \cap (A \cup B)$ का मिळते, हेही ओळखावे लागेल.

वरील सिद्धतेचे सर्व तपशील स्पष्ट केलेले आणखी एक रूप असे:
\begin{proof}{}
[संचांसाठीच्या $=$ च्या व्याख्येनुसार, $A \cap (A \cup B) = A$ दाखवण्यासाठी (a) $A \cap (A \cup B) \subseteq A$ आणि (b) $A \subseteq A \cap (A \cup B)$ दाखवावे लागते. (a): $\subseteq$ च्या व्याख्येनुसार $z \in A \cap (A \cup B)$ असेल, तर $z \in A$ असते हे दाखवावे लागते.] $z \in A \cap (A \cup B)$ असेल, तर $z \in A$ [$\cap$ च्या व्याख्येनुसार $z \in A \cap (A \cup B)$ हे $z \in A$ आणि $z \in A \cup B$ असते तेव्हाच आणि फक्त तेव्हाच सत्य असते]. म्हणून $A \cap (A \cup B) \subseteq A$. [(b): $\subseteq$ च्या व्याख्येनुसार $z \in A$ असेल, तर $z \in A \cap (A \cup B)$ असते हे दाखवावे लागते.] आता [(1)] $z \in A$ असे समजा. मग [(2)] $z \in A \cup B$ देखील [(1) वरून $z \in A$ किंवा $z \in B$; आणि $\cup$ च्या व्याख्येनुसार याचा अर्थ $z \in A \cup B$]. म्हणून $z \in A \cap (A \cup B)$ देखील [$\cap$ च्या व्याख्येनुसार $z \in A$, म्हणजे (1), आणि $z \in A \cup B$, म्हणजे (2), दोन्ही आवश्यक आहेत].
\end{proof}

\begin{prob}
$A \cup (A \cap B) = A$ ची पुढील सिद्धता विस्तृत करा. वापरलेल्या सर्व अनुमानाच्या नमुन्यांचा उल्लेख करा; प्रत्येक पायरी गृहीतकांवरून किंवा आधीच्या निष्कर्षांवरून का मिळते आणि कुठे कोणत्या व्याख्यांचा आधार घ्यावा लागतो, हे स्पष्ट करा.
\begin{proof}
$z \in A \cup (A \cap B)$ असेल, तर $z \in A$ किंवा $z \in A \cap B$. $z \in A \cap B$ असेल, तर $z \in A$. कोणताही $z \in A$ हा $\in A \cup (A \cap B)$ देखील असतो.
\end{proof}
\end{prob}












\section{मला जमत नाही!{}}\label{mth:prf:cnt:sec}

कधी ना कधी आपल्या सर्वांना हार मानावीशी वाटते. पण तुम्हाला हे \emph{जमू शकते}. शिक्षक, अध्यापन-सहायक आणि सहाध्यायी तुम्हाला मदत करू शकतात. त्यांच्याकडे मदत मागा!{} अडचणींचा मोठा पेच होऊ नये म्हणून आणि हार मानावीशी वाटल्यास काय करावे यासाठी काही सूचना पुढे दिल्या आहेत.

प्रश्न यशस्वीपणे सोडवता यावेत यासाठी पुढील गोष्टी करा:
\begin{enumerate}
\item \emph{शक्य तितक्या आधी सुरुवात करा.} सत्रभर आपण व्यस्त असतो आणि आपल्यापैकी अनेकांना काम पुढे ढकलण्याची सवय अडवते. गृहपाठाची सुरुवात लवकर करणे ही करता येण्याजोग्या सर्वात उपयुक्त गोष्टींपैकी एक आहे. त्यामुळे अडकलात तरी रडत बसण्यापलीकडे एखादा उपाय शोधण्यासाठी वेळ मिळतो.
\item \emph{सहाध्यायांशी बोला}. तुम्ही एकटे नाही. वर्गातील इतरांनाही अडचणी येत असतील; पण त्यांच्या अडचणी वेगळ्या असू शकतात. सहाध्यायांशी चर्चा केल्यामुळे प्रश्नाकडे पाहण्याची वेगळी दृष्टी मिळू शकते आणि कदाचित मार्ग सापडेल. अर्थात, त्यांचे उत्तर नुसते उतरवू नका: एखादी खूण मागा किंवा कुठे अडकलात हे समजावून पुढची पायरी विचारा. तुम्हाला समजले की त्यांनाही मदत करा. दुसऱ्याला पुढे जाण्यास मदत केल्याने आणि समजावून सांगितल्याने तुमची स्वतःची समजही सुधारते.
\item \emph{मदत मागा.} तुमच्यासाठी अनेक साधने उपलब्ध आहेत. शिक्षक आणि अध्यापन-सहायक तुमच्या मदतीसाठी आहेत आणि तुम्ही यशस्वी \emph{व्हावे अशी त्यांची इच्छा आहे}. प्रश्न सोडवण्यात मदत करणे आणि प्रक्रियेत नेमकी कुठे अडचण येते हे ओळखणे त्यांना शक्य असले पाहिजे.
\item \emph{थोडी विश्रांती घ्या.} अडकलात, तर त्या प्रश्नाकडे खूप वेळ पाहत बसल्यामुळे तसे झाले असण्याची \emph{शक्यता} आहे. थोडी विश्रांती घ्या, चहा प्या किंवा काही वेळ दुसरा प्रश्न सोडवा; मग ताज्या मनाने पहिल्या प्रश्नाकडे परत या. एखादी रात्र झोप घेऊन दुसऱ्या दिवशी पुन्हा पाहा.
\end{enumerate}

या उपायांसाठी सिद्धतेवर काम बराच वेळ आधी सुरू केलेले असणे आवश्यक आहे, हे लक्षात आले का? गृहपाठ देण्याच्या आदल्या दिवशी पहाटे दोन वाजता सिद्धता सुरू केली असेल, तर हे उपाय फार उपयोगी ठरणार नाहीत.

हे निराशाजनक वाटेल; पण सिद्धता करणे हे शेवटी समाधान देणारे आव्हान आहे. काही लोक हेच व्यवसाय म्हणून करतात; म्हणजे त्यात आनंद देणारे काहीतरी असले पाहिजे. जवळजवळ इतर सर्व गोष्टींप्रमाणे प्रश्न सोडवणे आणि सिद्धता करणे यांनाही सराव लागतो. काही सहाध्यायांना हे सोपे वाटत असेल; बहुधा त्यांनी आधीच बराच सराव केला असेल. फार लवकर हार मानू नका.

एखाद्या प्रश्नासाठी वेळ किंवा धीर संपला, तरी ठीक आहे. त्याचा अर्थ तुम्ही मूर्ख आहात किंवा तुम्हाला कधीच जमणार नाही असा नाही. शिक्षक किंवा दुसरा विद्यार्थी यांच्याकडून तो कसा सोडवतात ते समजून घ्या. कुठे चूक झाली किंवा कुठे अडकलात ते ओळखा, म्हणजे पुढे तशी अडचण आली तर ती टाळता येईल. मग उत्तर न पाहता पुन्हा सोडवण्याचा प्रयत्न करा. पुढच्या वेळी आधी सुरुवात करा आणि मदतही लवकर मागा.












\section{इतर साधने}\label{mth:prf:res:sec}

गणितात सिद्धता कशा कराव्यात याविषयी अनेक उपयुक्त पुस्तके आहेत. मूळ स्रोताने विशेषतः \emph{How to Read and do Proofs: An Introduction to Mathematical Thought Processes} (सोलो 2013) आणि \emph{How to Prove It: A Structured Approach} (व्हेलमन 2019) पाहण्याची सूचना केली आहे. \href{http://www.people.vcu.edu/~rhammack/BookOfProof/BookOfProof.pdf}{\emph{Book of Proof}} (हॅमॅक 2013) आणि \href{https://scholarworks.gvsu.edu/books/7/}{\emph{Mathematical Reasoning}} (सँडस्ट्रम 2019) ही सिद्धतांवरील पुस्तके मूळ स्रोतामध्ये इंटरनेटवर मुक्तपणे उपलब्ध म्हणून नमूद केलेली आहेत. तत्त्वज्ञानाच्या विद्यार्थ्यांसाठी \emph{More Precisely: The Math you need to do Philosophy} (स्टाइनहार्ट 2018) हे गणितीय विचाराचे उपयुक्त प्रास्ताविक पुस्तक ठरू शकते.

इंटरनेटवर सिद्धतांविषयी काही छोटी मार्गदर्शकेही आहेत. उदाहरणार्थ, \href{https://math.berkeley.edu/~hutching/teach/proofs.pdf}{“Introduction to Mathematical Arguments”} (हचिंग्ज 2003) आणि \href{https://eugeniacheng.com/wp-content/uploads/2017/02/cheng-proofguide.pdf}{“How to write proofs”} (चेंग 2004).

\subsection{प्रेरणादायी चित्रफिती}

गृहपाठ करण्याची प्रेरणाच उरली नाही असे वाटते का? मन खिन्न आहे का? या चित्रफिती मदत करू शकतील!

\begin{itemize}
\item \url{https://www.youtube.com/watch?v=ZXsQAXx_ao0}
\item \url{https://www.youtube.com/watch?v=BQ4yd2W50No}
\item \url{https://www.youtube.com/watch?v=StTqXEQ2l-Y}
\end{itemize}



